\begin{abstract}
Repairing a knowledge graph extracted from a document requires recovering
missing information while avoiding unsupported edits. We study a restricted
field-level setting with a declared schema and available source text.
The evaluated pipeline combines structural normalization, a diagnostic-context
language-model call, and deterministic candidate validation. On 225 held-out
document graphs containing 450 controlled defects, the archived pipeline
repairs 98.00\% of defects and achieves 99.17\% triple F1; a strong one-call
simple pipeline reaches 95.11\% repair and 96.34\% F1. A deterministic parser of the serialized source fields, however, reconstructs
all reference graphs without model calls; these results do not establish
superiority over source reconstruction. On 225 actual extraction
outputs, it raises F1 from 87.99\% to 92.48\%. Excluding ten malformed JSON
inputs, F1 improves from 92.09\% to 93.53\%, distinguishing factual repair
from format-failure recovery. A matched Gemma experiment holds instructions
and completion budget fixed, compares base, diagnostic, and SHACL contexts,
and evaluates filtering on exactly the same candidate responses. Diagnostic
context does not improve the matched baseline: controlled F1 after filtering
is 98.61\% with diagnostics and 99.62\% without them. Filtering itself yields
a smaller natural-input gain of 0.69 points in the diagnostic arm. We also
execute the earlier sequential optimizer on frozen proposals: it adds no
controlled repairs beyond preprocessing. Thus, the headline results support
the one-call pipeline and do not validate a neural or sequential optimization
mechanism. Evaluation uses structured silver references; independent human
adjudication remains pending. Code, prompts, predictions, and execution traces
are provided to make these boundaries auditable.
\keywords{Knowledge Graph Repair, Source Grounding, Constraint Validation, Mechanism Evaluation}
\end{abstract}
